\documentclass[reqno]{amsart}
\usepackage[T1]{fontenc}
\usepackage{amsmath,amssymb,amsthm}
\usepackage{xr-hyper}
\usepackage[hidelinks]{hyperref}
\externaldocument[ip-]{20261124T214600--the-isotropic-plane__workbook}
\externaldocument[pf-]{20261125T013700--pairings-on-a-formal-group-of-height-two__workbook}
\externaldocument[kl-]{20261125T009000--kronecker-limit-and-chowla-selberg__workbook}
\externaldocument[sm-]{20261125T000000--slope-method__workbook}
\externaldocument[ho-]{20261125T016500--hecke-operators-and-theta-lifts__workbook}
\externaldocument[c3-]{20261125T005000--galois-fourier-and-averaging__workbook}
\externaldocument[c32-]{20261126T030300--the-data__workbook}
\externaldocument[c35-]{20261126T030500--divisibility-of-the-composite-lift__workbook}

% Chapter number; \ref{chap} gives it.
\newcounter{chap}
% Chapter numbers, generated from 20261125T012000--grh-workbook__guide.org
% by scripts/make-chapters.py.  Do not edit: change the order of the
% index and rerun the script (./preview-site.sh and CI do this).
%
% Usage in a sheet, after \newcounter{chap}:
%   % Chapter numbers, generated from 20261125T012000--grh-workbook__guide.org
% by scripts/make-chapters.py.  Do not edit: change the order of the
% index and rerun the script (./preview-site.sh and CI do this).
%
% Usage in a sheet, after \newcounter{chap}:
%   % Chapter numbers, generated from 20261125T012000--grh-workbook__guide.org
% by scripts/make-chapters.py.  Do not edit: change the order of the
% index and rerun the script (./preview-site.sh and CI do this).
%
% Usage in a sheet, after \newcounter{chap}:
%   \input{chapters}\chapnumber{<sheet ID>}
\makeatletter
\providecommand\chapters@set[2]{\@namedef{chapters@#1}{#2}}
\chapters@set{20261125T001000}{1}% zeros-to-ramp-sums
\chapters@set{20261125T001100}{2}% inert-primes-and-estimate-b
\chapters@set{20261125T005000}{3}% galois-fourier-and-averaging
\chapters@set{20261125T005500}{4}% columns-from-a-sphere
\chapters@set{20261125T000000}{5}% slope-method
\chapters@set{20261126T030100}{6}% a-twist-for-each-prime
\chapters@set{20261125T003000}{7}% the-chain-and-the-square-root
\chapters@set{20261125T010000}{8}% rank-lemma-and-zero-estimate
\chapters@set{20261125T010500}{9}% poincare-series-and-near-orthogonality
\chapters@set{20261124T223000}{10}% elliptic-curves-for-the-proof
\chapters@set{20261125T009500}{11}% fock-model-and-theta-intertwiners
\chapters@set{20261124T223100}{12}% complex-multiplication-and-inert-reduction
\chapters@set{20261125T008000}{13}% one-cm-determinant
\chapters@set{20261124T214500}{14}% supersingular-theta-divisibility
\chapters@set{20261124T214600}{15}% the-isotropic-plane
\chapters@set{20261125T013700}{16}% pairings-on-a-formal-group-of-height-two
\chapters@set{20261125T013500}{17}% hecke-nondegeneracy-on-the-tree
\chapters@set{20261126T030200}{18}% the-bound-is-attained
\chapters@set{20261125T009000}{19}% kronecker-limit-and-chowla-selberg
\chapters@set{20261125T002000}{20}% definite-theta-lifts
\chapters@set{20261125T015500}{21}% siegel-weil-in-the-smallest-case
\chapters@set{20261125T002100}{22}% definite-theta-lifts-and-norms
\chapters@set{20261125T016500}{23}% hecke-operators-and-theta-lifts
\chapters@set{20261125T016700}{24}% divisibility-along-the-chain
\chapters@set{20261125T004000}{25}% denominators-and-sha
\chapters@set{20261125T004100}{26}% tate-twists-and-sha
\chapters@set{20261126T020430}{27}% the-minus-part-by-gauss-sums
\chapters@set{20261125T014100}{28}% congruence-modules
\chapters@set{20261125T014500}{29}% eisenstein-congruences-and-selmer-classes
\chapters@set{20261125T014700}{30}% pseudocharacters-and-the-length-bound
\chapters@set{20261125T011000}{31}% galois-representations-black-box
\chapters@set{20261126T030300}{32}% the-data
\chapters@set{20261126T020210}{33}% bessel-poincare-series-and-their-gram-operators
\chapters@set{20261126T030400}{34}% the-seesaw-identity
\chapters@set{20261126T030500}{35}% divisibility-of-the-composite-lift
\chapters@set{20261126T030600}{36}% theorem-a-short-forms
\chapters@set{20261126T030700}{37}% the-determinant-of-a-chain
\chapters@set{20261126T030800}{38}% the-cost-of-one-lift
\chapters@set{20261126T030900}{39}% denominators-of-projections-to-lifts
\chapters@set{20261126T031000}{40}% twists-with-few-congruences
\chapters@set{20261126T031100}{41}% theorem-c-and-the-main-inequality
\providecommand\chapnumber[1]{%
  \@ifundefined{chapters@#1}%
    {\PackageError{chapters}{Sheet #1 is not listed in the index}%
      {List it in 20261125T012000--grh-workbook__guide.org\MessageBreak
       and run scripts/make-chapters.py.}}%
    {\edef\thechap{\@nameuse{chapters@#1}}}}
\makeatother
\chapnumber{<sheet ID>}
\makeatletter
\providecommand\chapters@set[2]{\@namedef{chapters@#1}{#2}}
\chapters@set{20261125T001000}{1}% zeros-to-ramp-sums
\chapters@set{20261125T001100}{2}% inert-primes-and-estimate-b
\chapters@set{20261125T005000}{3}% galois-fourier-and-averaging
\chapters@set{20261125T005500}{4}% columns-from-a-sphere
\chapters@set{20261125T000000}{5}% slope-method
\chapters@set{20261126T030100}{6}% a-twist-for-each-prime
\chapters@set{20261125T003000}{7}% the-chain-and-the-square-root
\chapters@set{20261125T010000}{8}% rank-lemma-and-zero-estimate
\chapters@set{20261125T010500}{9}% poincare-series-and-near-orthogonality
\chapters@set{20261124T223000}{10}% elliptic-curves-for-the-proof
\chapters@set{20261125T009500}{11}% fock-model-and-theta-intertwiners
\chapters@set{20261124T223100}{12}% complex-multiplication-and-inert-reduction
\chapters@set{20261125T008000}{13}% one-cm-determinant
\chapters@set{20261124T214500}{14}% supersingular-theta-divisibility
\chapters@set{20261124T214600}{15}% the-isotropic-plane
\chapters@set{20261125T013700}{16}% pairings-on-a-formal-group-of-height-two
\chapters@set{20261125T013500}{17}% hecke-nondegeneracy-on-the-tree
\chapters@set{20261126T030200}{18}% the-bound-is-attained
\chapters@set{20261125T009000}{19}% kronecker-limit-and-chowla-selberg
\chapters@set{20261125T002000}{20}% definite-theta-lifts
\chapters@set{20261125T015500}{21}% siegel-weil-in-the-smallest-case
\chapters@set{20261125T002100}{22}% definite-theta-lifts-and-norms
\chapters@set{20261125T016500}{23}% hecke-operators-and-theta-lifts
\chapters@set{20261125T016700}{24}% divisibility-along-the-chain
\chapters@set{20261125T004000}{25}% denominators-and-sha
\chapters@set{20261125T004100}{26}% tate-twists-and-sha
\chapters@set{20261126T020430}{27}% the-minus-part-by-gauss-sums
\chapters@set{20261125T014100}{28}% congruence-modules
\chapters@set{20261125T014500}{29}% eisenstein-congruences-and-selmer-classes
\chapters@set{20261125T014700}{30}% pseudocharacters-and-the-length-bound
\chapters@set{20261125T011000}{31}% galois-representations-black-box
\chapters@set{20261126T030300}{32}% the-data
\chapters@set{20261126T020210}{33}% bessel-poincare-series-and-their-gram-operators
\chapters@set{20261126T030400}{34}% the-seesaw-identity
\chapters@set{20261126T030500}{35}% divisibility-of-the-composite-lift
\chapters@set{20261126T030600}{36}% theorem-a-short-forms
\chapters@set{20261126T030700}{37}% the-determinant-of-a-chain
\chapters@set{20261126T030800}{38}% the-cost-of-one-lift
\chapters@set{20261126T030900}{39}% denominators-of-projections-to-lifts
\chapters@set{20261126T031000}{40}% twists-with-few-congruences
\chapters@set{20261126T031100}{41}% theorem-c-and-the-main-inequality
\providecommand\chapnumber[1]{%
  \@ifundefined{chapters@#1}%
    {\PackageError{chapters}{Sheet #1 is not listed in the index}%
      {List it in 20261125T012000--grh-workbook__guide.org\MessageBreak
       and run scripts/make-chapters.py.}}%
    {\edef\thechap{\@nameuse{chapters@#1}}}}
\makeatother
\chapnumber{<sheet ID>}
\makeatletter
\providecommand\chapters@set[2]{\@namedef{chapters@#1}{#2}}
\chapters@set{20261125T001000}{1}% zeros-to-ramp-sums
\chapters@set{20261125T001100}{2}% inert-primes-and-estimate-b
\chapters@set{20261125T005000}{3}% galois-fourier-and-averaging
\chapters@set{20261125T005500}{4}% columns-from-a-sphere
\chapters@set{20261125T000000}{5}% slope-method
\chapters@set{20261126T030100}{6}% a-twist-for-each-prime
\chapters@set{20261125T003000}{7}% the-chain-and-the-square-root
\chapters@set{20261125T010000}{8}% rank-lemma-and-zero-estimate
\chapters@set{20261125T010500}{9}% poincare-series-and-near-orthogonality
\chapters@set{20261124T223000}{10}% elliptic-curves-for-the-proof
\chapters@set{20261125T009500}{11}% fock-model-and-theta-intertwiners
\chapters@set{20261124T223100}{12}% complex-multiplication-and-inert-reduction
\chapters@set{20261125T008000}{13}% one-cm-determinant
\chapters@set{20261124T214500}{14}% supersingular-theta-divisibility
\chapters@set{20261124T214600}{15}% the-isotropic-plane
\chapters@set{20261125T013700}{16}% pairings-on-a-formal-group-of-height-two
\chapters@set{20261125T013500}{17}% hecke-nondegeneracy-on-the-tree
\chapters@set{20261126T030200}{18}% the-bound-is-attained
\chapters@set{20261125T009000}{19}% kronecker-limit-and-chowla-selberg
\chapters@set{20261125T002000}{20}% definite-theta-lifts
\chapters@set{20261125T015500}{21}% siegel-weil-in-the-smallest-case
\chapters@set{20261125T002100}{22}% definite-theta-lifts-and-norms
\chapters@set{20261125T016500}{23}% hecke-operators-and-theta-lifts
\chapters@set{20261125T016700}{24}% divisibility-along-the-chain
\chapters@set{20261125T004000}{25}% denominators-and-sha
\chapters@set{20261125T004100}{26}% tate-twists-and-sha
\chapters@set{20261126T020430}{27}% the-minus-part-by-gauss-sums
\chapters@set{20261125T014100}{28}% congruence-modules
\chapters@set{20261125T014500}{29}% eisenstein-congruences-and-selmer-classes
\chapters@set{20261125T014700}{30}% pseudocharacters-and-the-length-bound
\chapters@set{20261125T011000}{31}% galois-representations-black-box
\chapters@set{20261126T030300}{32}% the-data
\chapters@set{20261126T020210}{33}% bessel-poincare-series-and-their-gram-operators
\chapters@set{20261126T030400}{34}% the-seesaw-identity
\chapters@set{20261126T030500}{35}% divisibility-of-the-composite-lift
\chapters@set{20261126T030600}{36}% theorem-a-short-forms
\chapters@set{20261126T030700}{37}% the-determinant-of-a-chain
\chapters@set{20261126T030800}{38}% the-cost-of-one-lift
\chapters@set{20261126T030900}{39}% denominators-of-projections-to-lifts
\chapters@set{20261126T031000}{40}% twists-with-few-congruences
\chapters@set{20261126T031100}{41}% theorem-c-and-the-main-inequality
\providecommand\chapnumber[1]{%
  \@ifundefined{chapters@#1}%
    {\PackageError{chapters}{Sheet #1 is not listed in the index}%
      {List it in 20261125T012000--grh-workbook__guide.org\MessageBreak
       and run scripts/make-chapters.py.}}%
    {\edef\thechap{\@nameuse{chapters@#1}}}}
\makeatother
\chapnumber{20261125T013500}
\renewcommand{\thesection}{\thechap.\arabic{section}}

\theoremstyle{definition}
\newtheorem{problem}{Problem}[section]
\theoremstyle{plain}
\newtheorem{theorem}{Theorem}[section]

\newcommand{\Z}{\mathbb{Z}}
\newcommand{\Q}{\mathbb{Q}}
\newcommand{\F}{\mathbb{F}}
\newcommand{\C}{\mathbb{C}}
\newcommand{\Sym}{\operatorname{Sym}}
\newcommand{\Hom}{\operatorname{Hom}}
\newcommand{\GL}{\mathrm{GL}}
\newcommand{\len}{\operatorname{length}}

\setlength{\emergencystretch}{2em}

\title{Hecke nondegeneracy on the tree}
\author{}
\date{November 25, 2026}

\begin{document}
\maketitle
\refstepcounter{chap}\label{chap}

\paragraph{Prerequisites.}
Modules over a discrete valuation ring and the elementary divisor theorem \cite[Ch.~III, \S7]{La}.  Unramified extensions of \(\Q_p\) and Teichm\"uller representatives.  The tree of \(\GL_2\) over a local field \cite[Ch.~II, \S1]{SeT}.  Legendre's formula \(v_p(m!) = (m - s_p(m))/(p-1)\), where \(s_p(m)\) is the sum of the base-\(p\) digits of \(m\).  From Chapter~\ref{sm-chap}, \emph{Size against divisibility}: Problem~\ref{sm-pr:fock}, and the colength of a submodule, defined before Theorem~\ref{sm-thm:colength}.  Section~\ref{sec:general} uses the notation of Chapter~\ref{c3-chap} for the places, and recalls the few quantities it needs from Chapter~\ref{c32-chap}.

In Chapter~\ref{sm-chap} an index was bounded by colengths of submodules (Theorem~\ref{sm-thm:colength}).  When a gain at a prime holds on a submodule spanned by images of a Hecke operator, how large is that colength?  For \(\GL_2\) and weight \(k\) we compute, at one fixed lattice, the image of the normalized sum of the values at the neighbouring lattices.  When \(Q > k\) its elementary divisors are expressed by binomial coefficients, and Kummer's theorem bounds them.

Why Hecke operators?  In the rewritten Lean proof \cite{GL} the maps of the chain are the lifts of the theta kernels \texttt{liftKernel} (\texttt{New/UpperBound/Objects/Columns.lean}), composed in \texttt{chain} (\texttt{New/UpperBound/Objects/LiftChains.lean}).  At an exceptional prime \(p < x_1\), in the sense of Section~\ref{sec:general}, the proof uses the image of a power of a product of scaled Hecke operators, modulo a power of \(p\).  Theorem~\ref{thm:general} at the end of this chapter bounds the colength needed for this step, and the computation at one vertex is a model for this bound.  Chapter~\ref{ho-chap}, \emph{Hecke operators and theta lifts}, relates a sum over the neighbours of a lattice on a class set to a Hecke operator on modular forms through the theta lift.  This chapter uses nothing from Chapter~\ref{ho-chap}.

Throughout, \(K\) is the unramified extension of \(\Q_p\) of degree \(f\), with ring of integers \(W\), residue field \(\F_Q\) where \(Q = p^f\), and valuation \(v\) normalized by \(v(p) = 1\).  We fix an integer \(k \ge 1\).

\section{The tree and the integral forms}\label{sec:tree}

A \emph{lattice} is a free \(W\)-submodule of rank \(2\) of \(K^2\).  For a lattice \(L\), its \emph{sources} are the lattices \(L''\) with \(pL \subset L'' \subset L\) and \(L/L'' \cong \F_Q\).  The vertices of the tree are the classes of lattices up to homothety, and the neighbours of the class of \(L\) are the classes of its sources \cite[Ch.~II, \S1]{SeT}.

\begin{problem}\label{pr:sources}
The sources of \(L\) correspond to the lines of \(L/pL \cong \F_Q^2\).  There are \(Q + 1\) of them, and their classes are pairwise distinct.
\end{problem}

Let \(V\) be the space of homogeneous polynomial functions of degree \(k\) on \(K^2\), on which \(\GL_2(K)\) acts by \((gF)(u) = F(g^{-1}u)\).  For a lattice \(L\) with dual \(L^\vee = \Hom_W(L, W)\), put
\[
  H(L) := \Sym^k(L^\vee) \subset V .
\]
If \(x, y\) is the basis of \(L^\vee\) dual to a basis of \(L\), then \(H(L) = \bigoplus_{j=0}^k W x^{k-j} y^j\).

\begin{problem}\label{pr:equivariance}
\(H(gL) = gH(L)\) for \(g \in \GL_2(K)\), and \(H(cL) = c^{-k} H(L)\) for \(c \in K^\times\).  If \(L'' \subset L\), then \(H(L) \subset H(L'')\).
\end{problem}

\begin{problem}\label{pr:one-edge}
Let \(L''\) be a source of \(L\).  There is a basis \(x, y\) of \(L^\vee\) such that \(L''^\vee = Wx + Wp^{-1}y\).  Hence the elementary divisors of \(H(L)\) in \(H(L'')\) are \(1, p, \dots, p^k\), we have \(p^k H(L'') \subset H(L)\), and \(\len_W H(L'')/H(L) = k(k+1)/2\).
\end{problem}

A \emph{form} is a function \(F\) from lattices to \(V\).  It is \emph{integral} if \(F(L) \in H(L)\) for every \(L\).  The Hecke operator and its normalization are
\[
  (TF)(L) := \sum_{L''} F(L''), \qquad Y := p^k T ,
\]
where the sum runs over the \(Q + 1\) sources of \(L\).  Let
\[
  S(L) := \sum_{L''} p^k H(L'') \subset H(L) .
\]

\begin{problem}\label{pr:hecke-integral}
\(Y\) maps integral forms to integral forms, and \(\{(YF)(L) : F \text{ integral}\} = S(L)\).
\end{problem}

\begin{problem}\label{pr:dual-side}
\(S(L) = \sum_M \Sym^k(M)\), the sum over the \(Q + 1\) lattices \(M\) with \(pL^\vee \subset M \subset L^\vee\) and \(L^\vee/M \cong \F_Q\).
\end{problem}

\begin{problem}\label{pr:small}
For \(Q = 2\) and \(k = 2\), \(S(L) = Wx^2 + 2Wxy + Wy^2\).  For \(k = 1\) and every \(Q\), \(S(L) = H(L)\).
\end{problem}

The \emph{loss} at \(L\) is \(\len_W H(L)/S(L)\).  It does not depend on \(L\), by Problem~\ref{pr:equivariance}.  By Problem~\ref{pr:one-edge} a single source costs \(k(k+1)/2\).  The question is how much the \(Q + 1\) sources save together.

\section{Elementary divisors}\label{sec:smith}

Let \(H\) be a free \(W\)-module of rank \(t\) and \(S \subset H\) a submodule of full rank.  By the elementary divisor theorem \cite[Ch.~III, \S7]{La} there is a basis \(h_1, \dots, h_t\) of \(H\) with \(S = \bigoplus_i p^{e_i} W h_i\), and the exponents \(e_1 \le \dots \le e_t\) are unique.  Then \(\len_W H/S = e_1 + \dots + e_t\).  For an integer \(D \ge 0\) and a real \(\varrho \in [0,1]\), we say that \(S\) is \emph{nondegenerate up to \(p^D\) on a share \(\varrho\)} when \(\#\{i : e_i \le D\} \ge \varrho t\).

\begin{problem}\label{pr:minors}
Let \(A\) be a matrix whose rows are the coordinates, in a basis of \(H\), of a finite generating set of \(S\).  For every integer \(1 \le l \le t\), the sum \(e_1 + \dots + e_l\) is the least valuation of an \(l \times l\) minor of \(A\).  In particular the \(e_i\) do not change under a finite unramified extension of \(W\).
\end{problem}

\begin{problem}\label{pr:good-count}
\(\#\{i : e_i \le D\}\) is the largest \(\nu\) for which there are \(s_1, \dots, s_\nu \in S\) and \(\varphi_1, \dots, \varphi_\nu \in \Hom_W(H, W)\) with \(\varphi_a(s_b) = p^D \delta_{ab}\).
\end{problem}

The vectors \(s_b\) are \emph{witnesses} and the functionals \(\varphi_a\) are \emph{tests}.  The next problem gives a criterion on linear combinations of the tests.

\begin{problem}\label{pr:detect}
Let \(\Lambda\) be a \(W\)-module and \(\varphi_1, \dots, \varphi_M \colon \Lambda \to W\) linear maps.  Suppose that for every \(c \in W^M\) with some \(c_i\) a unit there is \(G \in \Lambda\) with \(v\bigl(\sum_i c_i \varphi_i(G)\bigr) \le D\).  Then the image of \(G \mapsto (\varphi_i(G))_i\) contains \(p^D W^M\).
\end{problem}

\begin{problem}\label{pr:trunc}
Let \(A, A'\) be \(t \times t\) matrices over \(W\) with \(A \equiv A' \pmod{p^T}\) for an integer \(T \ge 1\), and let \(e_i\) and \(e'_i\) be the exponents of their row spans in \(W^t\) (where \(e_i = \infty\) for the zero elementary divisors).  Then \(\min(e_i, T) = \min(e'_i, T)\) for every \(i\).
\end{problem}

Problem~\ref{pr:trunc} lets one compute exponents below \(T\) modulo \(p^T\); Section~\ref{sec:colength} measures a loss in this truncated way.

\section{Kummer's theorem}\label{sec:kummer}

Valuations of binomial coefficients enter the formula for the exponents.

\begin{problem}\label{pr:kummer}
For \(0 \le j \le n\), \(v_p\binom{n}{j} = \bigl(s_p(j) + s_p(n-j) - s_p(n)\bigr)/(p-1)\), and this is the number of carries in the addition of \(j\) and \(n - j\) in base \(p\).
\end{problem}

This is due to Kummer \cite{Ku}; see \cite{Gra} for a survey.

\begin{problem}\label{pr:kummer-bound}
For \(1 \le j \le n\), \(v_p\binom{n}{j} \le \lfloor \log_p n \rfloor - v_p(j)\).  In particular \(p^{v_p\binom nj} \le n\).
\end{problem}

\begin{problem}\label{pr:kummer-cases}
Every \(\binom{Q-1}{j}\) with \(0 \le j \le Q-1\) is prime to \(p\).  For \(0 < j < p^m\), \(v_p\binom{p^m}{j} = m - v_p(j)\), so Problem~\ref{pr:kummer-bound} is sharp.
\end{problem}

\section{The Fock form}\label{sec:fock}

For \(F = \sum_m F_m x^m\) and \(G = \sum_m G_m x^m\) in \(K[x_1, \dots, x_{n'}]\), with multi-indices \(m\) and \(m! = \prod_s m_s!\), let
\[
  \langle F, G \rangle := \sum_m F_m G_m\, m! ,
  \qquad\text{that is,}\qquad
  \langle x^m, x^{m'} \rangle = m!\, \delta_{m m'} .
\]
In one variable over \(\C\), with complex conjugation on the coefficients of \(F\), this is the Gaussian inner product of Problem~\ref{sm-pr:fock}; it is the inner product of the Fock model \cite{Ba}.

\begin{problem}\label{pr:fock-adj}
\(\langle x_s F, G \rangle = \langle F, \partial_s G \rangle\) for every \(s\).  Hence \(\langle F, G \rangle = \bigl(F(\partial) G\bigr)(0)\), where \(F(\partial)\) is \(F\) with each \(x_s\) replaced by \(\partial/\partial x_s\).
\end{problem}

\begin{problem}\label{pr:fock-power}
For \(F \in K[x, y]\) homogeneous of degree \(k\) and \(a, b \in K\), \(\langle (ax + by)^k, F \rangle = k!\, F(a, b)\).
\end{problem}

Let \(H = \bigoplus_{j=0}^k W x^{k-j} y^j\), and let \(H^\sharp := \{G \in K[x,y] \text{ homogeneous of degree } k : \langle G, F \rangle \in k!\, W \text{ for all } F \in H\}\).

\begin{problem}\label{pr:fock-dual}
\(H^\sharp = \bigoplus_{j=0}^k \binom kj W x^{k-j} y^j\), and \(p^{\lfloor \log_p k \rfloor} H \subset H^\sharp \subset H\).
\end{problem}

So \(\langle \cdot, \cdot \rangle / k!\) is a perfect pairing between \(H^\sharp\) and \(H\).  By Kummer's theorem, each elementary-divisor exponent of \(H^\sharp \subseteq H\) is at most \(\lfloor \log_p k \rfloor\), so the inclusion is nondegenerate up to \(p^{\lfloor \log_p k \rfloor}\) on all of the rank.

\section{The span at a vertex}\label{sec:vertex}

Fix \(L\) and a basis \(x, y\) of \(L^\vee\), so that \(H(L) = H\) and \(H(L)^\sharp = H^\sharp\) in the notation of Section~\ref{sec:fock}.  The lattices \(M\) of Problem~\ref{pr:dual-side} are \(M_\alpha := W(x + \alpha y) + pL^\vee\) for \(\alpha \in W\), which depends only on \(\alpha\) modulo \(p\), and \(M_\infty := Wy + pL^\vee\).

\begin{problem}\label{pr:powers}
(i) If \(Q > k\), the \(W\)-span of \(\{\lambda^k : \lambda \in L^\vee\}\) is \(H(L)^\sharp\).  (ii) If \(Q = 2\) and \(k = 3\), this span has colength \(1\) in \(H(L)^\sharp = H(L)\).
\end{problem}

\begin{problem}\label{pr:span-i}
\(\Sym^k(M_\alpha)\) is spanned by the \((x + \alpha y)^{k-i}(py)^i\) with \(0 \le i \le k\).  If \(Q > k\), then for each \(i\) the \(W\)-span of \(\{(x + \alpha y)^{k-i} y^i : \alpha \in W\}\) is \(\bigoplus_{j=i}^k \binom{k-i}{j-i} W x^{k-j} y^j\).  For the remaining line, \(\Sym^k(M_\infty) = \bigoplus_{j=0}^k p^{k-j} W x^{k-j} y^j\).
\end{problem}

Problems~\ref{pr:dual-side} and~\ref{pr:span-i} give the following theorem.  Problem~\ref{pr:kummer-bound} gives its last two sentences.

\begin{theorem}\label{thm:vertex}
Suppose \(Q > k\).  Then \(S(L) = \bigoplus_{j=0}^k p^{\varepsilon_j} W x^{k-j} y^j\), where the exponent \(\varepsilon_j\) of the monomial \(x^{k-j}y^j\) is
\[
  \varepsilon_j = \min\Bigl(k - j,\ \min_{0 \le i \le j} \Bigl(i + v_p\binom{k-i}{j-i}\Bigr)\Bigr) .
\]
In particular \(\varepsilon_j \le \min\bigl(j,\, k-j,\, v_p\binom kj\bigr) \le \lfloor \log_p k \rfloor\).  So \(p^{\lfloor \log_p k \rfloor} H(L) \subset S(L)\), and the loss is at most \(\sum_j v_p\binom kj\).
\end{theorem}

When \(Q > k\), the inclusion \(S(L) \subseteq H(L)\) is therefore nondegenerate up to \(p^{\lfloor \log_p k \rfloor}\) on all of the rank.  The exponents can be strictly smaller than the valuations \(v_p\binom kj\).

\begin{problem}\label{pr:check}
Assume \(Q > k\).  Then \(\varepsilon_j = \varepsilon_{k-j}\).  For \(k = Q - 1\) the loss is \(0\).  For \(k = p < Q\) the loss is \(p - 1\).  For \(p = 3\), \(k = 6\) and \(Q \ge 9\) the \(\varepsilon_j\), in the order \(j = 0, \dots, k\), are \(0, 1, 1, 0, 1, 1, 0\).  For \(p = 2\), \(k = 8\) and \(Q \ge 16\) they are \(0, 1, 1, 1, 1, 1, 1, 1, 0\), so the loss is \(7\), although \(\sum_j v_2\binom 8j = 17\).
\end{problem}

What if there are few edges?

\begin{problem}\label{pr:few-edges}
If \(Q \le k\), then at most \(Q + 1\) of the exponents of \(S(L)\) in \(H(L)\) vanish, and the loss is at least \(k - Q\).  More generally, for each integer \(u \ge 1\) at most \((Q+1)u\) of the exponents are less than \(u\).
\end{problem}

So the largest exponent is at least \((k+1)/(Q+1) - 1\), and for fixed \(Q\) the loss grows quadratically in \(k\), since at least \(k + 1 - (Q+1)u\) exponents are \(\ge u\) for each \(u\).  For \(Q = 2\) and \(k = 15\), a Smith form computation gives the exponents \(0, 0, 0, 1, 1, 1, 2, 2, 2, 3, 3, 3, 4, 4, 4, 5\), a loss of \(35\).  For \(Q = 3\) and \(k = 15\) it gives \(0, 0, 0, 0, 1, 1, 1, 1, 2, 2, 2, 2, 4, 4, 4, 4\).  The formula of Theorem~\ref{thm:vertex} assumes \(Q > k\).

When \(Q > k\), every exponent is at most \(\lfloor \log_p k \rfloor\).  Conversely, if every exponent is at most \(D\), then \(Q + 1 \ge (k+1)/(D+1)\).

\section{Onto up to colength}\label{sec:colength}

At an exceptional prime \(p < x_1\), the rewritten Lean proof \cite{GL} bounds the colength of a Hecke image after adding a power of \(p\) times the module of integral source forms.  Problem~\ref{pr:trunc} illustrates this truncation over \(W\).  The colength of Chapter~\ref{sm-chap} makes sense over any commutative ring.  Let \(O\) be a commutative ring, \(\mathcal V\) an \(O\)-module, \(\mathcal V_1, \mathcal V_2 \subseteq \mathcal V\) submodules and \(\zeta \in O\).  We say that \emph{the colength of \(\mathcal V_2\) in \(\mathcal V_1\) divides \(\zeta\)} if for every \(t \ge 0\) and every family \(v_1, \dots, v_t\) of elements of \(\mathcal V_1\) there is a \(t \times t\) matrix \(C\) over \(O\) with \(\det C \mid \zeta\) and \(\sum_j C_{ij} v_j \in \mathcal V_2\) for every \(i\).  No inclusion between \(\mathcal V_1\) and \(\mathcal V_2\) is assumed.  For \(p \in O\), an integer \(u \ge 0\), an endomorphism \(U\) of \(\mathcal V\) and a real number \(c\), we say that \emph{\(U\) is onto \(\mathcal V_1\) modulo \(p^u\) up to colength \(c\)} if \(U\mathcal V_1 \subseteq \mathcal V_1\) and there is an integer \(l \ge 0\) with \(l \le c\) such that the colength of \(U\mathcal V_1 + p^u \mathcal V_1\) in \(\mathcal V_1\) divides \(p^l\).

\begin{problem}\label{pr:colength}
Let \(\mathcal V_1\) be a free \(W\)-module of finite rank, \(\mathcal V_2 \subseteq \mathcal V_1\) a submodule, possibly with \(\mathcal V_1/\mathcal V_2\) of infinite length, and \(l \ge 0\) an integer.  The colength of \(\mathcal V_2\) in \(\mathcal V_1\) divides \(p^l\) if and only if \(\len_W \mathcal V_1/\mathcal V_2 \le l\).
\end{problem}

\begin{problem}\label{pr:onto-vertex}
Suppose \(Q > k\), and let \(U\) be an endomorphism of \(H(L)\) with \(U H(L) = S(L)\), for example the diagonal one that multiplies \(x^{k-j}y^j\) by \(p^{\varepsilon_j}\).  Then \(U\) is onto \(H(L)\) modulo \(p^u\) up to colength \(c\) if and only if \(c \ge \sum_{j=0}^k \min(\varepsilon_j, u)\).  For \(p = 3\), \(k = 6\), \(Q \ge 9\) and every \(u \ge 1\) the least such \(c\) is \(4\).
\end{problem}

\begin{problem}\label{pr:onto-prod}
Let \(\mathcal V_1\) be a free \(W\)-module of finite rank.  If \(U_1\) and \(U_2\) are onto \(\mathcal V_1\) modulo \(p^u\) up to colengths \(c_1\) and \(c_2\), then \(U_1 U_2\) is onto \(\mathcal V_1\) modulo \(p^u\) up to colength \(c_1 + c_2\).
\end{problem}

The rewritten Lean proof proves the rule of Problem~\ref{pr:onto-prod} over any commutative ring, for a product of any finite number of endomorphisms.

\section{Exceptional primes and the theorem on elementary divisors}\label{sec:general}

Where does the rewritten proof use a product of scaled Hecke operators?  We use the notation of Chapter~\ref{c3-chap}: a prime \(d\) and a modulus \(q_0\), divisible by \(18480\), with \(q_0 \mid 2(d+1)\) and \(q_0 \nmid d+1\); the quadratic character \(\chi\) of \(\Q(\sqrt{-d})\); the infinite places of the field \(E\), identified with \((\Z/q_1)^\times\) for a multiple \(q_1\) of \(q_0\) that has the same prime factors as \(q_0\); and the \(a\)-places, the classes \(\tau \equiv 1 \pmod{q_0}\).  For the application to the chain, fix an admissible datum as in Definition~\ref{c32-def:data}; write \(m\) for its number of lifts, \(x < x_1 < b\) for its scales, and \(a = xb\).  The gain asked of the chain at a prime \(p\) with \(\chi(p) = -1\) is a sum of terms indexed by places \(\sigma\) and odd \(e \ge 1\), each involving the place \(-p^e\sigma\).  For a prime \(p\) with \(\chi(p) \ne 1\), let \(I_p\) be the set of odd \(e \ge 1\) if \(\chi(p) = -1\), and of all \(e \ge 1\) if \(p = d\).  A \emph{collision} is a pair \((\tau, e)\) of an \(a\)-place \(\tau\) and an \(e \in I_p\) for which \(-p^e\tau\) is again an \(a\)-place.  There are none at \(p = d\), since \(d \equiv q_0/2 - 1 \pmod{q_0}\) gives \(-d^e \equiv 1 - q_0/2\) or \(-1 \pmod{q_0}\).  Since \(d \equiv -1\) modulo \(q_0/2\), which is divisible by \(8 \cdot 3 \cdot 5 \cdot 7 \cdot 11\), quadratic reciprocity gives \(\chi = 1\) at every prime factor of \(q_0\); so a prime \(p\) with \(\chi(p) = -1\) is prime to \(q_0\), and we write \(o_0(p)\) for the order of \(-p\) in \((\Z/q_0)^\times\).  Let \(R_p\) be the least integer \(j \ge 0\) with \(m \le p^j\).  Following the rewritten Lean proof, \(p\) is \emph{exceptional} if \(\chi(p) = -1\), \(p \not\equiv \pm 1 \pmod m\), \(o_0(p)\) is odd and \(o_0(p) \le R_p\).

\begin{problem}\label{pr:collision}
Let \(p\) be a prime with \(\chi(p) = -1\), and \(e \ge 1\) odd.  Then \((\tau, e)\) is a collision for one, equivalently every, \(a\)-place \(\tau\) if and only if \(o_0(p) \mid e\).  So collisions exist if and only if \(o_0(p)\) is odd, and if \(o_0(p)\) is odd and \(o_0(p) > R_p\), every collision has \(p^e \ge pm\).
\end{problem}

The rewritten Lean proof charges the collisions to an allowed loss.  For each prime \(p\), fix the embedding \(\iota_p \colon \C \to \C_p\) used to test integrality, let \(\mathcal F\) be the complex space of source forms, and \(\mathcal F_{\mathrm{int}}\) its module of integral forms.  Write \(f_E = [E:\Q(\sqrt{-d})]\) and \(\mathrm{loss}(p) = f_E(\log b)^2 a/p\), multiplied by \(m\) if \(p \equiv \pm 1 \pmod m\).  At a prime \(p < x_1\) with \(\chi(p) \ne 1\) that is not exceptional, the sum of \(ma/p^e\) over all collisions is at most \(\mathrm{loss}(p)\), and the proof takes the whole module \(\mathcal F_{\mathrm{int}}\), with colength exponent \(0\).  If \(\chi(p) = 1\) or \(p \ge x_1\), the gain is \(0\).

At an exceptional prime \(p < x_1\) the proof chooses a \emph{full mark} \(\Lambda\): a set of primes of \(E\), each induced by an \(a\)-place, disjoint from its conjugate, and containing, for every \(a\)-place, the prime it induces or the conjugate of that prime.  For \(w \in \Lambda\) let \(A_w\) be \(a\) times the number of \(a\)-places inducing \(w\), and \(T_w\) the Hecke operator at \(w\) of the double coset of \(\operatorname{diag}(\varpi_w, \dots, \varpi_w, 1, \dots, 1)\), with \(R_p\) entries equal to a uniformizer \(\varpi_w\).  Let \(U\) be the ordered product of the operators \(p^{A_w}T_w\).  The proof takes \(U^m\mathcal F_{\mathrm{int}} + p^{j_p}\mathcal F_{\mathrm{int}}\), with \(j_p = \lceil 2mf_Ea/p \rceil\).  The operator \(U^m\) is onto \(\mathcal F_{\mathrm{int}}\) modulo \(p^{j_p}\) up to colength \(\mathrm{loss}(p)\dim_{\C}\mathcal F\); this follows by adding the bounds for the factors of \(m\) copies of the product, by the rule of Problem~\ref{pr:onto-prod}, and estimating the sum.  The bound for each factor uses the following theorem.

Let \(N \ge 1\), let \(W'\) be a commutative local ring with maximal ideal \(\mathfrak m\), and let \(O\) be a valuation domain of characteristic \(0\).  In the setting of the theorem, a \emph{Frobenius cycle} of length \(N\) from \(W'\) to \(O\) consists of a prime \(p\), a power \(P \ge 5\) of \(p\), ring homomorphisms \(\iota_s \colon W' \to O\) for \(s \in \Z/N\), a parity, even or odd, for each \(s\), and an element \(\varpi \in O\), such that \(\varpi^2 = p\); \(p\) divides \(\iota_{s+1}(w) - \iota_s(w)^P\) in \(O\) for all \(s\) and all \(w \in W'\); the parities of \(s\) and \(s+1\) differ for every \(s\); and \(W'/\mathfrak m\) has \(q := P^N\) elements.  Fix integers \(\kappa \ge 1\) (the rank), \(\mu \ge 0\) (the degree in each letter \(s\)) and \(r \ge 0\).  Let \(\mathcal B\) be the set of families \(\beta = (\beta_s)_{s \in \Z/N}\) of maps \(\beta_s \colon \{0, \dots, \kappa-1\} \to \Z_{\ge 0}\) with \(\sum_i \beta_s(i) = \mu\) for each \(s\).  For \(\gamma \in \GL_\kappa(W')\) put \(\gamma_s = \iota_s(\gamma)\) if \(s\) is even and \(\gamma_s = \iota_s((\gamma^{-1})^{\mathsf T})\) if \(s\) is odd, and let \(\rho(\gamma)\) be the \(\mathcal B \times \mathcal B\) matrix whose \((\beta, \beta')\) entry is the product over \(s\) of the coefficients of \(z^{\beta_s}\) in \(\prod_i \bigl(\sum_j (\gamma_s)_{ji} z_j\bigr)^{\beta'_s(i)}\).  The \emph{depth} \(\mathrm{dp}(\beta)\) is the sum of \(\beta_s(i)\) over the pairs \((i, s)\) with \(s\) even and \(i \ge r\), or \(s\) odd and \(i < r\); and \(\Delta\) is the diagonal matrix with entries \(p^{\mathrm{dp}(\beta)}\).  Let \(\mathcal K_r \subseteq \GL_\kappa(W')\) consist of the matrices whose \((i, j)\) entry lies in \(\mathfrak m\) whenever \(i < r \le j\), and let \(\mathcal C = \GL_\kappa(W')/\mathcal K_r\), the set of left cosets.

For the theorem take finite sets \(X\) and \(Z\).  \emph{Edge coordinates} on \((X, Z)\) are maps \(\pi \colon X \times \mathcal C \to X\), \(\gamma \colon X \times \mathcal C \to \GL_\kappa(W')\), \(\gamma' \colon \mathcal C \to \GL_\kappa(W')\) and \(\sigma \colon X \times \mathcal C \to \GL_Z(O)\), such that \(\gamma'(\omega)^{-1} \mathcal K_r = \omega\) for every \(\omega \in \mathcal C\), and \(\eta = \eta'\) whenever \(\eta, \eta' \in X \times \mathcal C\) satisfy \(\pi(\eta) = \pi(\eta')\) and \(\gamma(\eta)^{\mathsf T} (\gamma(\eta')^{\mathsf T})^{-1} \in \mathcal K_r\).  For a finite set \(I\) and a bijection \(J = (J_1, J_2) \colon I \to \mathcal B \times Z\), the \emph{edge Hecke operator} \(\mathcal A\) is the \(O\)-linear endomorphism of \((O^I)^X\) given by
\[
  (\mathcal A F)(\xi)_i = \sum_{\omega \in \mathcal C} \sum_{i' \in I} \sigma(\xi, \omega)_{J_2 i,\, J_2 i'}\, \bigl(\rho(\gamma(\xi, \omega))\, \Delta\, \rho(\gamma'(\omega))\bigr)_{J_1 i',\, J_1 i}\; F(\pi(\xi, \omega))_{i'} .
\]
For \(\kappa = 2\) and \(r = 1\) the cosets \(\mathcal C\) correspond to the lines of the residue plane, as the sources of Problem~\ref{pr:sources} do.

\begin{theorem}\label{thm:general}
Let a Frobenius cycle of length \(N \ge 1\) from a commutative local ring \(W'\) to a valuation domain \(O\) of characteristic \(0\) be given, with prime \(p\), power \(P\) and \(q = P^N\) as above, and suppose that \(p\) is not a unit in \(O\).  Let \(\kappa, \mu, r\) be integers and \(\theta\) a real number with, for the natural logarithm,
\[
  0 < r, \quad 2r \le \kappa, \quad 0 < \theta \le 1, \quad \mu^4 \le q, \quad \Bigl(\frac{N \kappa \log q}{\theta}\Bigr)^5 P^{3r} \le \mu .
\]
Let \(X\), \(Z\) and \(I\) be finite sets (\(\mathcal C\) is finite), \(J \colon I \to \mathcal B \times Z\) a bijection, edge coordinates on \((X, Z)\) with edge Hecke operator \(\mathcal A\), \(\mathcal M\) an \(O\)-module, \(\psi \colon (O^I)^X \to \mathcal M\) an \(O\)-linear map and \(U\) an endomorphism of \(\mathcal M\) with \(U \psi = \psi \mathcal A\).  Then for every integer \(u \ge 0\), \(U\) is onto \(\mathcal L = \psi\bigl((O^I)^X\bigr)\) modulo \(p^u\) up to colength
\[
  \bigl(3\,c_{\mathrm{div}} + \theta u\bigr)\, \#X\, \#I, \qquad c_{\mathrm{div}} = N \mu \Bigl(\frac{1}{P^{2r}} + \frac{r^2}{\kappa P}\Bigr) .
\]
\end{theorem}

The bound has a part \(\theta u \,\#X\,\#I\), a share \(\theta\) of the bound \(u\,\#X\,\#I\) that every endomorphism keeping \(\mathcal L\) satisfies, as \(\mathcal L\) is generated by \(\#X\,\#I\) elements, and a part \(3c_{\mathrm{div}}\#X\,\#I\) that does not depend on \(u\).  It bounds a colength, not the individual elementary divisors.  Theorem~\ref{thm:vertex} is its model, not a case of it:

\begin{problem}\label{pr:cycle}
A Frobenius cycle has even length \(N\).  Under the hypotheses on \(\kappa, \mu, r, \theta\) of Theorem~\ref{thm:general}, \(\mu > P^{3r}\) and \(N > 12r\).  For \(N = 14\), \(\kappa = 2\), \(r = 1\), \(\theta = 1/2\), \(P = 5^{100}\) and \(\mu = 5^{350}\) these hypotheses hold.
\end{problem}

\paragraph{In the Lean formalization.}
In the rewritten Lean proof \cite{GL}, readable edition, the notions of Section~\ref{sec:colength} are \texttt{ColengthDvd} and \texttt{IsOntoUpToColength}, with the product rule \texttt{IsOntoUpToColength.prod} (\texttt{New/Reduction/Lattices/Colength.lean}); exceptional primes are \texttt{IsExceptional} (\texttt{New/UpperBound/Size/ShortColumns.lean}); and Theorem~\ref{thm:general} is \texttt{FrobeniusCycle.elementary\_divisors} (\texttt{New/UpperBound/ElementaryDivisors/ElementaryDivisorTheorem.lean}).  The statements on the other primes are \texttt{mul\_collisionSum\_le\_loss} (\texttt{New/UpperBound/Divisibility/JetExponent.lean}), \texttt{gain\_eq\_zero} (\texttt{New/UpperBound/Size/Estimates.lean}) and the proof of \texttt{exists\_colengthDvd\_gaugeLe\_compositeLift}; the bound for \(U^m\) is \texttt{IsMarked.isOntoUpToColength\_prod\_scaledHeckeOperator\_pow}, both in \texttt{New/UpperBound/Divisibility/LocalDivisibility.lean}.  The bound for one factor, \texttt{IsMarked.isOntoUpToColength\_scaledHeckeOperator} (\texttt{New/UpperBound/Divisibility/MarkOperator.lean}), applies the theorem through \texttt{Columns.isOntoUpToColength\_scaledHeckeOperator\_of\_frobeniusCycle} in \texttt{New/UpperBound/Divisibility/HeckeEdges.lean}, with \(P = p^{o_0(p)}\), \(\mu = a\), \(r = R_p\) and \(\theta = 1/(m^2 f_E)\), the letters being the \(a\)-places that induce the prime of the operator or its conjugate.  The other objects of Section~\ref{sec:general} are, in \texttt{New/Reduction/Data/}, \texttt{Parameters}, \texttt{Parameters.q}\(_1\) and \texttt{Parameters.placeMulEquiv} (\texttt{FieldsAndPlaces.lean}) and \texttt{Data.aPlaces} (\texttt{AdmissibleData.lean}); in \texttt{ShortColumns.lean}, \texttt{localPrimeSum}, \texttt{gain}, \texttt{loss}, \texttt{IsCollision}, \texttt{collisionSum}, \texttt{collisionExponent}, \texttt{o}\(_0\) and \texttt{ceilLogLifts}; in \texttt{New/UpperBound/Divisibility/MarkedPrimes.lean}, \texttt{IsFullMark}, \texttt{integralSourceForms}, \texttt{heckeScalingExponent} and \texttt{IsMarked.scaledHeckeOperator}, with \texttt{Data.SourceForm} in \texttt{New/UpperBound/Objects/Columns.lean}, \texttt{heckeSplit} in \texttt{Lib/Weil/Unramified/SplitHecke.lean} and \texttt{basicDoubleCoset} in \texttt{Lib/Adelic/Unitary/Defs.lean}; in \texttt{New/UpperBound/ElementaryDivisors/}, \texttt{FrobeniusCycle}, \texttt{Exponents}, \texttt{letterRep}, \texttt{depth} (\texttt{TwistingData.lean}), \texttt{ParahoricCosets} (\texttt{Edges.lean}), \texttt{EdgeCoordinates}, \texttt{edgeHeckeOperator\_apply}, \texttt{IsLargeDegree} (\texttt{ElementaryDivisorTheorem.lean}) and \texttt{divisorConstant} (\texttt{Depth.lean}).

Chapters~\ref{ip-chap} and~\ref{pf-chap} study the local gain, and this chapter the colength at an exceptional prime below \(x_1\); Chapter~\ref{c35-chap} combines the local estimates for the composite lift.  The next part, from Chapter~\ref{kl-chap} on, turns to periods, masses and norms of theta lifts.

\section*{Hints}

\subsection*{\S\ref{sec:tree}}
\begin{description}
\item[\ref{pr:sources}] \(L'' \mapsto L''/pL\).  For the last statement, compare the indices in \(L\) of two homothetic sources.
\item[\ref{pr:equivariance}] \((gL)^\vee = (g^{-1})^{\mathsf T} L^\vee\) under the pairing.  An inclusion of lattices reverses on duals.
\item[\ref{pr:one-edge}] Choose a basis \(u_1, u_2\) of \(L\) with \(L'' = Wu_1 + Wpu_2\), by the elementary divisor theorem.
\item[\ref{pr:hecke-integral}] For the image, take \(F\) supported on the sources of \(L\).
\item[\ref{pr:dual-side}] Dualize the two inclusions defining a source.
\item[\ref{pr:small}] List the three lattices \(M\) and expand.
\end{description}

\subsection*{\S\ref{sec:smith}}
\begin{description}
\item[\ref{pr:minors}] Multiplying \(A\) on either side by an invertible matrix does not change the ideal generated by its \(l \times l\) minors.
\item[\ref{pr:good-count}] To construct a family, use the dual basis of \((h_i)\).  To bound the size of any such family, reduce the witness and test maps modulo \(p^{D+1}\).
\item[\ref{pr:detect}] The image is a submodule of \(W^M\); compare it with \(p^D W^M\).
\item[\ref{pr:trunc}] Compare the cokernels of \(A\) and \(A'\) after tensoring with \(W/p^T W\).
\end{description}

\subsection*{\S\ref{sec:kummer}}
\begin{description}
\item[\ref{pr:kummer}] Legendre's formula, and each carry lowers the digit sum by \(p - 1\).
\item[\ref{pr:kummer-bound}] Between which digit positions can a carry occur?
\item[\ref{pr:kummer-cases}] Problem~\ref{pr:kummer}.
\end{description}

\subsection*{\S\ref{sec:fock}}
\begin{description}
\item[\ref{pr:fock-adj}] Check it on monomials.  Then induct on the degree of \(F\).
\item[\ref{pr:fock-power}] Problem~\ref{pr:fock-adj}, or expand both sides.
\item[\ref{pr:fock-dual}] The monomials are orthogonal.  For the inclusions, Problem~\ref{pr:kummer-bound}.
\end{description}

\subsection*{\S\ref{sec:vertex}}
\begin{description}
\item[\ref{pr:powers}] Use the Vandermonde matrix of Teichm\"uller representatives.  For \(Q = 2\), look for a cubic form that takes values in \(2\Z_2\) on \(\Z_2^2\).
\item[\ref{pr:span-i}] The same Vandermonde argument, applied to the vectors \(\bigl(\binom{k-i}{j-i} \alpha^{j-i}\bigr)_j\).
\item[\ref{pr:check}] The symmetry \(x \leftrightarrow y\) of \(H(L)\) permutes the lattices \(M\).  Then Problem~\ref{pr:kummer-cases}.
\item[\ref{pr:few-edges}] Reduce modulo \(p\).  For the general statement, use Problem~\ref{pr:one-edge} modulo \(p^u\).
\end{description}

\subsection*{\S\ref{sec:colength}}
\begin{description}
\item[\ref{pr:colength}] For one direction take \(v_1, \dots, v_t\) a basis of \(\mathcal V_1\).  For the other, the vectors \(c \in W^t\) with \(\sum_j c_j v_j \in \mathcal V_2\) form a submodule whose quotient embeds in \(\mathcal V_1/\mathcal V_2\).
\item[\ref{pr:onto-vertex}] Theorem~\ref{thm:vertex} and Problem~\ref{pr:colength}; then Problem~\ref{pr:check}.
\item[\ref{pr:onto-prod}] \(U_1\) maps \(\mathcal V_1/(U_2\mathcal V_1 + p^u\mathcal V_1)\) onto \((U_1\mathcal V_1 + p^u\mathcal V_1)/(U_1U_2\mathcal V_1 + p^u\mathcal V_1)\).
\end{description}

\subsection*{\S\ref{sec:general}}
\begin{description}
\item[\ref{pr:collision}] For odd \(e\), \(-p^e = (-p)^e\).
\item[\ref{pr:cycle}] Go once around \(\Z/N\).  For the size, \(N\kappa\log q/\theta > 1\); for the example, \(\log 5 < 2\).
\end{description}

\begin{thebibliography}{SeT}
\bibitem[Ba]{Ba} V.~Bargmann, \emph{On a Hilbert space of analytic functions and an associated integral transform, Part~I}, Comm.\ Pure Appl.\ Math.\ \textbf{14} (1961), 187--214.
\bibitem[GL]{GL} L.~Sernau, \emph{GRH for Dirichlet \(L\)-functions, in Lean~4}, a shorter rewriting of the Lean proof, text by Claude, \url{https://github.com/lsernau/grh-lean}, commit \texttt{72fead1}, readable edition \texttt{grh/}.
\bibitem[Gra]{Gra} A.~Granville, \emph{Arithmetic properties of binomial coefficients I: Binomial coefficients modulo prime powers}, in: Organic Mathematics (Burnaby, 1995), CMS Conf.\ Proc.\ \textbf{20}, Amer.\ Math.\ Soc., 1997, 253--276.
\bibitem[Ku]{Ku} E.~E.~Kummer, \emph{\"Uber die Erg\"anzungss\"atze zu den allgemeinen Reciprocit\"atsgesetzen}, J.\ reine angew.\ Math.\ \textbf{44} (1852), 93--146.
\bibitem[La]{La} S.~Lang, \emph{Algebra}, rev.\ 3rd ed., Grad.\ Texts in Math.\ 211, Springer, 2002.
\bibitem[SeT]{SeT} J.-P.~Serre, \emph{Trees}, Springer, 1980.
\end{thebibliography}

\end{document}
